Super — ich habe deine **finale PDF** gründlich gecheckt und liefere dir ein klares, priorisiertes Ergebnis. Ich gehe **Step by Step** vor:

---

# 0) Kurzfazit

* **Go mit Auflagen**: Inhaltlich stark und gut strukturiert; **einige Blocker** müssen vor Upload behoben werden: **offene Zitate `[...]`**, **kaputtes PlantUML-Listing**, **mehrere Sprach-/Stil-Fehler**, **einige DOI/Referenz-Inkonsistenzen**, **s/w-Tauglichkeit & Fremdmaterial-Hinweise**. 

---

# 1) Springer/LNCS-Compliance (Formales)

**Erfüllt / unkritisch**

* **Umfang** passend (27 Seiten); LNCS-Layout; nummerischer Zitierstil. 
* **Abstract/Keywords** jetzt LNCS-tauglich (Problem → Contribution → Nutzen). 

**Prüf-/Fix-Punkte**

1. **s/w-Druck**: Mehrere Screenshots/Abbildungen sind farb- bzw. Graustufen-sensibel (z. B. Fig. 9/10/12). **Hatching/Markerformen/Linienstile** ergänzen; wenn möglich **Vektorgrafiken statt Screenshots**. (Springer mailte: Druck in s/w.) 
2. **Fremdmaterial/Permissions**:

   * Fig. 4 weist explizit eine Website-Quelle aus → **nachzeichnen** (eigene Tabelle) **oder** schriftliche Erlaubnis beilegen. 
   * Fig. 9 (ProBlind-Screenshot) & Fig. 12 (LTspice) → **Logo/Marke minimieren** bzw. **als eigene Vektor-Beispiele** nachbauen. 
3. **Running heads**: Bitte sicherstellen, dass `\titlerunning{…}` / `\authorrunning{…}` **kurz** gesetzt sind (im PDF nicht sichtbar, im .tex prüfen).
4. **Footnote/URL-Form**: In *Used Terms* steht ein nackter Link („1 https://…“) → als **echte Fußnote** setzen: `\footnote{\url{…}}`. 
5. **Unaufgelöste Literaturhinweise**: In mehreren *Related Work*-Abschnitten stehen **`[?]`-Platzhalter** → **Blocker** (Bib/Keys nachziehen, **bbl** neu bauen). Beispiele in **Pattern 1** und **Pattern 5**. 

---

# 2) Sprache & Rechtschreibung (konkrete Funde → Fix)

> Format: **Original** → **Vorschlag** (Begründung)

* **„are hardly or not at all suitable…“** → **„are often inaccessible…“** (idiomatischer, prägnant). 
* **„user defined LATEXBeamer template“** → **„user-defined \LaTeX{} Beamer template“** (Bindestrich, Marke, \LaTeX). 
* **„rendered documents e.g., as PDF documents“** → **„rendered outputs (e.g., PDF)“** (Kürzung, Stil, Komma nach *e.g.*). 
* **„Markdown … Hence, it is easier to parse …“**: Komma nach „Hence,“ ist da; **okay**. (Stelle insgesamt stimmig.) 
* **„Listing 1.1: … Fig 4“** → **„Fig.~4“** (Punkt & geschütztes Leerzeichen; LNCS-Konvention). 
* **„assistance tools“** (mehrfach) → **„assistive tools/technologies“** (Fachsprache). 
* **„… and 6 shows the generated diagram.“** → **„… and Fig.~\ref{fig:GeneratedUML} shows …“** (kein nacktes „6“, saubere Referenz). 
* **Satzanfang klein**: *„in [5] …“* → **„In [5] …“** (Großschreibung am Satzanfang). 
* **Zeitreihen-Stil**: *„… not well suited to representing and convey the time based aspect …“* →
  **„… not well suited to representing the time-based aspect …“** **oder**
  **„… not well suited to represent and convey the time-based aspects …“** (Parallelismus, Bindestrich). 
* **Kleinigkeiten**: vereinzelte fehlende Leerzeichen nach *e.g.,* („e.g., a …“ ist an vielen Stellen korrekt; bitte projektweit kurz prüfen). 

**Kaputtes Listing (Blocker)**

* **PlantUML-Code** enthält zerrissene Tokens: „**+ f i r s tname : S t r ing**“, „**Person −̂− Student**“ etc. → Copy/Paste-Artefakte; **Diagramm wird so nicht reproduzierbar**. **Ersetze** durch eine saubere Version (siehe Patch unten). 

---

# 3) Inhalt & Argumentation

**Sehr gut**

* **Problem → Gap → Contribution** ist im Abstract und in der Einleitung klar strukturiert: Lehr-Kontexte (UML, KV-Map, Schaltungen) → Barriere → Pattern-Lösungen. 
* **Fünf Patterns** decken den didaktischen Raum ab (Text-first, Diagramme, Grafiken/Taktik, Sonifikation, Generierte Erklärungen). **Beispiele und Toolketten** sind praxisnah. 

**Verbesserung (schnell machbar)**

* **Pattern-Namen konsistent** halten: In Text/Grafik teils „**Core of Content**“, teils *„Transformable Core Content“* in früheren Fassungen. Entscheide dich (ich empfehle **Transformable Core Content** wegen des Transformations-Kerns). 
* **Pattern 1 „Related Work“**: inhaltlich stimmig, aber aktuell **`[?]`** statt Zitationen → Bib-Keys eintragen (Austin/Seo/Cervone-Sorge/Armano). **Blocker**. 
* **Pattern 5 „Related Work“**: dito — drei `[…]`-Platzhalter durch **Leotta/Edwards/Jung** ersetzen. **Blocker**. 
* **Pattern 4**: Sehr gute Beispiele (RC-Schaltung, Highcharts). Ergänze **Mini-Checkliste** (Playback-Speed, Keyboard-Control, „numbers on demand“, Legend Voiceover) als 3–4 Bullets → steigert didaktischen Nutzen. 
* **Screenshots** (Fig. 9/10/12): Kurz **didaktische Pointe** in Captions ergänzen (Was/Warum/Wie-zugänglich), und s/w-Hinweis (Hatching/Marker). 

---

# 4) Literatur (Kohärenz & DOIs)

**Unaufgelöste Zitate**

* In **Pattern 1 & 5** stehen **`[?]`** → **zwingend** Bib-Keys einsetzen und **bbl** neu bauen. 

**Inkonsistenzen (bitte korrigieren)**

* **Stefik, 2007**: Eintrag sagt *„Proceedings of the 12th International ACM SIGACCESS Conference“*, DOI ist aber **IEEE/ICPC** (*10.1109/ICPC…*) → **Venue/DOI widersprechen sich**; den **korrekten DOI und Venue** nachtragen. 
* **Supalo, 2012**: Im Text steht *Journal of Chemical Education*, der DOI in der Referenz ist **JSESD** (*10.14448/jsesd…*) → **Journal/DOI mismatch**; richtigstellen. 
* **Brock/Jouffrais (2015)**: Venue-Bezeichnung (HCI-Konferenz) vs. **ACM-DOI** → bitte Titel/Venue/DOI konsistent machen. 

> Du hast bereits eine **aktualisierte `.bib`** von mir bekommen; nimm diese als Basis und gleiche die o. g. Einträge noch einmal gezielt ab.

---

# 5) Konkrete Patches (ready-to-paste)

**A) PlantUML (Listing 1.2) ersetzen**

```text
@startuml
class Person {
  + firstname : String
  + surname   : String
}
class Student {
  + matriculationNumber : int
  + degreeProgram       : String
}
Person ^-- Student
class Lecturer {
  + title      : String
  + department : String
}
Person ^-- Lecturer
class Book {
  + title : String
  + ISBN  : String
}
Student "1" o-- "n" Book
@enduml
```

→ Danach **Fig. 6** neu erzeugen; im Text „… and **Fig.~\ref{fig:GeneratedUML}** shows …“ einsetzen. 

**B) Footnote korrigieren (Used Terms)**

```latex
... visual impairment.\footnote{\url{https://lifeafterblindness.com/understanding-the-differences-between-blindness-and-visual-impairment/}}
```



**C) Formulierungen (Auswahl)**

* *„user defined LATEXBeamer“* → *„user-defined \LaTeX{} Beamer“* 
* *„assistance tools“* → *„assistive tools“* (projektweit) 
* *„Listing 1.1 … Fig 4“* → *„Fig.~4“* 
* *„not well suited to representing and convey …“* → *„… representing the time-based aspect …“* 

**D) Offene Zitate ersetzen (Blocker)**

* **Pattern 1**: `\cite{Austin2024PreTeXt,Seo2019RMarkdown,CervoneSorge2019W4A,Armano2018LaTeXPDF}` statt `[…]`. 
* **Pattern 5**: `\cite{Leotta2023AutoCaptioning,Edwards2023AltTextProcess,Jung2021AltText}` statt `[…]`. 

---

# 6) Priorität (kurz)

**Muss sofort**

1. **`[?]`-Zitate** fixen (P1 & P5) → **Bib/Build** neu. 
2. **PlantUML-Listing + Fig. 6** korrigieren. 
3. **Fremdmaterial/Permissions** (Fig. 4/9/12) klären oder nachzeichnen. 
4. **s/w-Tauglichkeit** aller Figuren (Hatching/Marker). 
5. **Referenz-Inkonsistenzen** (Stefik 2007, Supalo 2012, Brock 2015) bereinigen. 

**Sollte**

* Sprachfeinschliff (assistive vs. assistance; *e.g.,*-Komma; Fig.~) & Captions mit **didaktischer Pointe + A11y-Hinweis**. 

**Kann**

* Mini-Checklisten am Ende von Pattern 3/4 (Designregeln Taktik / Sonifikation-Usability). 

---

Wenn du willst, baue ich dir jetzt eine **Search&Replace-Patchliste** (Overleaf-freundlich) und markiere die **konkreten Zeilen** in deinen `.tex`-Dateien, damit du die 5 Blocker in <15 Min. schließen kannst.
